% =====================================================================
\section{Related work}\label{sec:related}
% =====================================================================
This section places the article among studies of workload balance in storage and station assignment, storage decisions under uncertain or changing demand, robust optimisation and out-of-sample evaluation, and it states what the article takes from each and where it differs from the closest work.

\subsection{Workload balance in pick-and-pass and zone storage}\label{subsec:rw_balance}
How work is spread over stations matters for the performance of a pick-and-pass line. \citet{yu2008pickandpass} approximate such systems as a queueing network to estimate order lead time and station utilisation, and use the approximation to evaluate storage methods at the stations. Balance has long been an aim of storage assignment itself. \citet{jane2000storage} assigns products to the zones of a relay picking line from historical orders so that pickers carry nearly equal loads, and adds heuristics that adjust locations when order volume fluctuates. \citet{tarczynski2023} formulates a mixed-integer model for pick-and-pass systems that minimises the maximum expected zone workload, and notes that assignment based on historical data carries risk, because even perfectly balanced zones can receive different workloads for a particular set of orders. \citet{kim2022dynamic} balance a weighted sum of picks and visits across the zones of a progressive bypass zone picking system and extend the model to reassignment under limited relocation capacity. \citet{zhang2024joint} impose a workload-balance constraint on picking aisles in a robotic mobile fulfilment system and report that robot movement distance increases as that constraint becomes more stringent. The companion \citep{belul_companion} minimises station visits subject to hard per-station budgets, each set at 110\% of the complete legacy load the station carried.

The closest work is \citet{winkelmann2025integrated}, who assign products to the stations and shelves of a pick-and-pass loop in an e-grocery fulfilment centre subject to constraints that bound each station's relative deviation in daily picks from the average over all stations from above and below. An assignment based on average demand exceeded a 3\% station deviation on individual days of the week against a 1\% threshold, and day-of-week-specific constraints brought the deviation below 1\%. Under simulated week-to-week variation with a coefficient of variation of 0.05, the maximum station deviation exceeded 4\% for the average-based assignment and still reached 2.46\% for the day-of-week-aware one. The authors suggest robust or stochastic optimisation for random variation as future work. Their study establishes that balance constraints set on average demand can be exceeded when demand varies. This article differs in four respects (Table~\ref{tab:closest}). Its targets are the frozen historical shares of each station under the incumbent layout, which differ from station to station, rather than the average over stations. Its band is absolute in share units rather than a relative deviation from that average. It protects the band against later demand with a reserved margin or with historical block scenarios rather than with day-of-week-specific constraints. It scores layouts on later observed orders of the same stream rather than on simulated weeks.

\begin{table}[htbp]
\centering
\caption{Closest studies to this article: how station or zone workload balance is stated, protected against demand change, and evaluated.}
\label{tab:closest}
\footnotesize
\setlength{\tabcolsep}{3pt}
\begin{tabular}{@{}p{2.2cm}p{2.3cm}p{2.7cm}p{2.6cm}p{2.5cm}@{}}
\toprule
Study & Setting & Balance rule & Protection against demand change & Evaluation \\
\midrule
\citet{tarczynski2023} & Pick-and-pass zones & Minimise the maximum expected zone workload & None; expected demand & Exemplary instances \\
\citet{kim2022dynamic} & Progressive bypass zone picking & Weighted picks and visits balanced across zones & Capacity-limited reassignment & Simulation with industry orders \\
\citet{winkelmann2025integrated} & Pick-and-pass loop & Upper and lower bounds on the relative deviation of daily picks from the average over all stations & Day-of-week-specific constraints & Simulated week-to-week variation \\
\citet{dundar2025robust} & Forward picking area & None & Budgeted robustness on correlation coefficients in the objective & Generated correlations, 5 to 12 SKUs \\
Companion \citep{belul_companion} & Pick-and-pass stations & One-sided budget, 110\% of legacy load & None & In sample; visits also on unseen dated weeks \\
This article & Pick-and-pass stations & Two-sided absolute band around each station's historical share & Reserved margin; historical block scenarios with activation & Later observed orders at one held-out origin \\
\bottomrule
\end{tabular}
\end{table}

\subsection{Storage decisions under uncertain or changing demand}\label{subsec:rw_uncertain}
Two lines of work respond to demand that is uncertain or changes over time. The first protects a plan. \citet{ang2012robust} model multiperiod storage and retrieval in a unit-load warehouse facing uncertain demand with partially characterised distributions, and minimise the worst-case expected total travel with a linear decision rule. \citet{dundar2025robust} derives a budgeted robust counterpart for correlated SKU-to-location assignment in which a budget limits how many uncertain correlation coefficients may deviate; the robustness applies to objective coefficients rather than constraints, no workload balance is modelled, and the author tests instances of 5, 10 and 12 SKUs and reports that larger instances become computationally intractable. Neither protects a workload restriction.

The second line changes the plan. \citet{christofides1973rearrangement} sequence the moves that rearrange a warehouse when relative demand changes, \citet{pazour2015reshuffling} optimise reshuffling by sequential moves, \citet{kim2022dynamic} limit reassignment to a relocation capacity, and \citet{wu2026reassignment} analyses when item storage should be reassigned as the demand pattern evolves in robotic mobile fulfilment. \citet{winkelmann2025integrated} observe that rearranging products is often not practical for short-term variation, whereas seasonal or long-term demand changes typically require it. The article takes no position on when to re-slot. It asks whether one re-slotting keeps the band on the orders that follow it.

Few of the reviewed studies score a storage allocation built from historical orders on orders that came later. \citet{mirzaei2021ica} build a cluster-based allocation from historical turnover and affinity and validate it on a real dataset in which 70\% of the orders, chosen at random, built the allocation and the rest were held out, a random rather than temporal split. \citet{waubertdepuiseau2022dslap} train a reinforcement-learning storage policy on one year of data and evaluate it on the following two months, a single temporal split. In the studies we reviewed, the companion and this article included, we found none that evaluates a storage decision across several successive origins.

\subsection{Robust optimisation ideas used here}\label{subsec:rw_robust}
The share rows of Section~\ref{sec:model} follow the set-inclusive, worst-case form of \citet{soyster1973convex}: a row must hold for every product mix in a declared set. \citet{bental1999robust} showed that robust counterparts of uncertain linear programmes can remain tractable, conic quadratic for ellipsoidal sets. Here the set is a polytope with finitely many relevant endpoints, so the counterpart consists of finitely many linear rows. \citet{bertsimas2003robust,bertsimas2004price} control conservatism with a per-constraint protection level that limits how many coefficients may deviate from their nominal values. The linear or integer structure is kept, and probabilistic bounds on constraint violation are attached. The reserve used in this article is different: a fixed fraction $\lambda$ of the band, subtracted on every share row, with no probabilistic bound attached. \citet{bertsimas2018datadriven} build uncertainty sets from data through statistical hypothesis tests, and their robust solutions carry finite-sample probabilistic guarantees. The scenario set used here is instead the convex hull of historical block mixes, extended by activation mass on historically inactive products; it is not a hypothesis-test set. Sampled-constraint methods relate feasibility to samples, through the exact feasibility of randomised solutions of uncertain convex programmes \citep{campi2008exact} or through sample approximations of chance constraints \citep{luedtke2008sample}.

The historical blocks of this article resemble sampled scenarios but do not meet the assumptions behind those guarantees. The guarantees of budgeted protection, of sample approximation and of distributionally robust data-driven decisions rest on independence or random-sampling assumptions \citep{bertsimas2004price,luedtke2008sample,vanparys2021data}. Consecutive blocks of one order stream are not randomly sampled, and their independence is not established. The model therefore claims only feasibility conditional on the future mix lying in the declared set, and the article implements none of the cited methods and tests none of their guarantees.

\subsection{Out-of-sample evaluation of optimised decisions}\label{subsec:rw_oos}
An optimised decision should not be judged on the data it was fitted to. \citet{smith2006optimizer} show that when alternatives are ranked by estimated value and the best is chosen, its true value should be expected to fall below its estimate even when the estimates are unbiased. \citet{vanparys2021data} define out-of-sample disappointment as the probability that the actual expected cost of a data-driven decision exceeds its predicted cost. The same reasoning, applied here to band compliance rather than cost, is why a layout optimised on history is scored on orders it was not fitted to. \citet{tashman2000outofsample} reviews out-of-sample tests of forecasting accuracy. \citet{cerqueira2020evaluating} note, following \citet{tashman2000outofsample}, that applying hold-out evaluation over multiple test periods gives more reliable estimates than a single partition, and their experiments on non-stationary series support this. By that standard the held-out evidence of this article is a single test period at one origin, preceded by an exploratory origin whose data it contains; the absence of further origins is a limitation of the study (Section~\ref{subsec:limitations}), not a gap it fills.

\subsection{What is standard and what this article adds}\label{subsec:rw_position}
Several components of the study are standard. Reserving part of a constraint is related to the protection levels of budgeted robust optimisation \citep{bertsimas2004price}, with the fixed-fraction difference stated in Section~\ref{subsec:rw_robust}. A linear row that holds at every point of a finite hull holds at its vertices, so worst-case rows over a scenario hull are a finite set of rows. Rescaling rows with integer data to integer counts with rounded thresholds uses only integrality. Scoring a decision on a later hold-out window is established in forecasting evaluation \citep{tashman2000outofsample,cerqueira2020evaluating}, and the companion already does so on unseen dated weeks, although few of the reviewed storage studies do (Section~\ref{subsec:rw_uncertain}). None of these is claimed as new.

Relative to the companion, the article replaces one-sided workload budgets by two-sided share rows under an explicit information contract. It compares a reserved margin with historical block scenarios in a predeclared factorial scored on later observed orders, and it reports diagnostics showing that the model's own feasibility does not explain the observed compliance. Relative to \citet{winkelmann2025integrated}, it differs in the respects stated in Section~\ref{subsec:rw_balance} and Table~\ref{tab:closest}, at the cost of evidence from one warehouse and one held-out origin. Among the studies reviewed, we found none that protects station workload-share constraints in pick-and-pass storage with a reserved margin or with scenarios built from historical orders and then scores the layouts on later observed orders; this is the article's reading of the reviewed set, not a claim about all published work.
